\chapter{Symmetry and the harmonic theta mechanism}\label{ch:theta}
\section{An involution that preserves the weight}
The coordinate labels are residue labels, so negation alone does not preserve the standard shell. The correct involution is
\begin{equation}\label{eq:j}
J(v)_i=-v_{(-i\bmod5)}.
\end{equation}
It preserves the residues, sum and norm. Its coordinate permutation is $(1\ 4)(2\ 3)$, which is even. Negation contributes $(-1)^{10}=1$ to the degree-$10$ Vandermonde. Hence
\[
W(Jv)=W(v).
\]
Write $s_m(v)=\sum_i v_i^m$. Then $s_m(Jv)=(-1)^ms_m(v)$. In particular, every odd homogeneous symmetric polynomial $Q$ satisfies $Q(Jv)=-Q(v)$. Pairing by $J$ gives
\begin{equation}\label{eq:oddnull}
\sum_{v\in\Vn}W(v)Q(v)=0.
\end{equation}
At a fixed point of $J$, such a $Q$ is zero, so the argument covers both orbit sizes. Odd corrections alone do not change the sum of an entire $J$-orbit. Even corrections will be the useful source of transfers between orbits.

\section{Why the Vandermonde is harmonic}
Work on the four-dimensional Euclidean hyperplane
\[
H=\{v\in\R^5:\sum_i v_i=0\}.
\]
The Laplacian on $H$ commutes with coordinate permutations. Since $D$ is alternating, its Laplacian is alternating and has degree $8$. Every alternating polynomial on $H$ vanishes when $v_i=v_j$. The ten distinct linear factors $v_i-v_j$ therefore divide it. A degree-$8$ polynomial cannot be divisible by their degree-$10$ product unless it is zero. Thus
\[
\Delta_H D=0.
\]
This degree argument is useful because it does not require expanding the ten-factor product.

More generally, if $Q$ is symmetric and $DQ$ is harmonic homogeneous of degree $10+d$, its weighted shell theta series is
\begin{equation}\label{eq:thetaq}
\Theta_Q(q)=\sum_{n\ge1}\left(\sum_{v\in\Vn}W(v)Q(v)\right)q^n.
\end{equation}
The harmonic theta theorem initially supplies a vector-valued modular form. Obtaining a scalar full-level cusp form requires the following lattice specialization, rather than an unexplained assertion that every shell theta series has level one.

\section{The regular discriminant cosets}
Let $A_4=H\cap\Z^5$ and $L=\sqrt5 A_4$. Its discriminant module is
\[
L^*/L\cong A_4^*/5A_4,
\qquad |L^*/L|=5^5=3125.
\]
A representative of $L^*$ can be written
\[
x=\frac1{\sqrt5}(a_i+t/5)_{i=0}^4,
\qquad0\le t<5,\qquad\sum_i a_i=-t.
\]
If two $a_i$ have the same residue modulo $5$, the transposition of those coordinates fixes the discriminant coset. An alternating theta component on that coset must equal its own negative and is zero.

If all five residues are distinct, they are $0,1,2,3,4$ in some order. Their sum is zero modulo $5$, so $\sum_i a_i=-t$ forces $t=0$. These regular cosets form one free $S_5$-orbit of size $120$. The sign representation occurs exactly once in the corresponding permutation module. It does not occur in an irregular orbit, since an odd stabilizer kills the alternating part. Thus the alternating theta vector lies on a one-dimensional sign line.

For a regular representative, $\sum_i a_i^2$ is divisible by $5$, since the squared residues sum to $30$. It is even, since $a_i^2\equiv a_i\pmod2$ and $\sum_i a_i=0$. The quadratic exponent $\sum_i a_i^2/10$ is integral. The Weil operator $T$ therefore acts as $1$ on the sign line. The Weil action commutes with $S_5$, so this line is invariant. The rank is even, and the representation descends to $SL_2(\Z)$. On the line the relation $(ST)^3=S^2$ reduces to $S^3=S^2$, giving $S=1$.

The full-level scalar conclusion has now been justified by the discriminant calculation. The ordinary residue coset has $x=v/\sqrt5$, and its exponent $\|x\|^2/2$ is exactly $\|v\|^2/10=n$. Scaling the homogeneous polynomial changes an overall scalar, not the modular weight or the zero conclusion.

\section{The modularity contract and the cusp condition}
The authenticated harmonic theta theorem states that a homogeneous harmonic polynomial of degree $\nu$ on a rank-$r$ positive definite even lattice gives a vector-valued modular form of weight $\nu+r/2$ for the Weil representation \cite{IW}. Here $r=4$ and $\nu=10+d$. The project-specific sign-line argument above therefore gives
\begin{equation}\label{eq:thetalevelone}
\Theta_Q\in S_{12+d}(SL_2(\Z))
\end{equation}
for the even-degree harmonic construction. The zero vector contributes no constant term. The full modular group has one cusp, so the form is cuspidal. This is the exact external-to-internal dependency: the general theorem supplies vector modularity, and the explicit coset argument supplies the scalar full-level specialization.

The convergence needed for the analytic theta function is harmless here. The lattice is positive definite, the polynomial has finite degree and Gaussian decay dominates polynomial growth on each compact subset of the upper half-plane. Termwise differentiation is therefore valid. The corresponding formal Fourier coefficients are finite shell sums.

\chapter{Universal null polynomials}\label{ch:nulls}
\section{The correction that does not change any coefficient}
A universally null polynomial is a symmetric polynomial $Q$ with
\begin{equation}\label{eq:universalnull}
\sum_{v\in\Vn}W(v)Q(v)=0\quad\text{for every }n\ge1.
\end{equation}
This is much stronger than being null on one selected shell. If a finite family satisfies \eqref{eq:universalnull}, any rational linear combination of it does too. Thus
\[
F(v)=1+\sum_j a_jQ_j(v)
\]
has the exact property
\begin{equation}\label{eq:correctedtotal}
\sum_{v\in\Vn}W(v)F(v)=\tau(n)
\end{equation}
at every index, after the coefficients $a_j$ have been fixed. The choice of coefficients may depend on a target shell. The null identity does not thereby become merely a shell-specific identity.

The retained construction obtains \eqref{eq:universalnull} by combining harmonicity with a finite vanishing certificate and the valence formula. This is one place where finite checks support a universal implication for a stated mathematical reason.

\section{A low-degree limit}
Every alternating polynomial of total degree at most $13$ is $D$ times a symmetric quotient of degree at most $3$. On $s_1=0$, such a quotient is a linear combination of $1,s_2,s_3$. On a shell, $s_2=10n$, while the $s_3$ contribution vanishes under $J$. Its weighted sum is thus $(a+10bn)\tau(n)$ for constants $a,b$.

This low-degree space cannot extract a new independent arithmetic residual. It is a scoped limitation on alternating polynomial refinements through total degree $13$. At quotient degree $4$, $s_4$ appears, and harmonic theta forms of higher weights introduce new moment identities. These will still require care: new cusp coefficients do not automatically give a nonzero lower bound for the original one.

\section{An exact operator on symmetric power sums}
Define $\mathcal L$ by
\[
\Delta_H(DQ)=D\mathcal LQ.
\]
The Vandermonde is harmonic. Its logarithmic derivative and the product rule give, with $s_0=5$ and $s_1=0$,
\begin{equation}\label{eq:lsm}
\mathcal Ls_m
=m\sum_{a=0}^{m-2}s_as_{m-2-a}
-\frac{m(m-1)}5s_{m-2},\qquad m\ge2.
\end{equation}
The chain-rule metric is
\begin{equation}\label{eq:gradient}
\nabla_Hs_m\cdot\nabla_Hs_n
=mn\left(s_{m+n-2}-\frac{s_{m-1}s_{n-1}}5\right).
\end{equation}
For a polynomial in $s_2,s_3,s_4,s_5$, these two formulas determine $\mathcal LQ$ by the ordinary second-order product rule. Higher power sums are reduced through Newton's identities for five variables.

The apparent denominators from differentiating $D$ cancel pairwise. The quotient
\[
\frac{v_i^{m-1}-v_j^{m-1}}{v_i-v_j}
\]
is a polynomial. Thus the operator really has rational polynomial coefficients on the symmetric ring. For example, \eqref{eq:lsm} gives $\mathcal Ls_2=48$. This agrees with the direct radial identity $\Delta_H(s_2D)=48D$.

At each even invariant degree $d$, one lists monomials $s_2^as_3^bs_4^cs_5^e$ with $2a+3b+4c+5e=d$. Writing an unknown rational linear combination, the equation $\mathcal LQ=0$ is a finite linear system. The initial theta coefficients supply further rational linear equations. No floating-point approximation is needed to define this kernel.

\section{Why a finite initial segment suffices}
The level-one valence formula implies that a nonzero holomorphic modular form of weight $k$ has order at the cusp at most $k/12$ \cite{Stein}. A cusp form whose coefficients of $q^1,\ldots,q^{\lfloor k/12\rfloor}$ all vanish has order at least $\lfloor k/12\rfloor+1>k/12$. It must be zero.

For a harmonic $DQ$ of invariant degree $d$, the theta weight is $k=12+d$. Thus the constructor imposes
\[
[q^m]\Theta_Q=0,\qquad1\le m\le\left\lfloor\frac{12+d}{12}\right\rfloor.
\]
Harmonicity supplies membership in the specified modular-form space, and valence supplies the implication from this finite certificate to all coefficients. Without harmonicity and the modularity theorem, testing the same initial segment would have no universal force.

The distinction between the three degree conventions is essential. The symmetric correction has invariant degree $d$. The alternating product $DQ$ has total degree $10+d$. Its theta form has modular weight $12+d$. A claim about corrections through degree $20$ means the first of these unless stated otherwise.

\section{The first explicit corrections}
The construction gives
\begin{align}\label{eq:q8}
Q_8={}&6273s_2^4-33750s_2^2s_4-56800s_2s_3^2\notag\\
&+69600s_3s_5+45000s_4^2,
\end{align}
and
\begin{align}\label{eq:q10}
Q_{10}={}&2091s_2^5-11250s_2^3s_4-145800s_2^2s_3^2\notag\\
&+409600s_2s_3s_5+15000s_2s_4^2\notag\\
&-30000s_3^2s_4-264000s_5^2.
\end{align}
Both $DQ_8$ and $DQ_{10}$ are harmonic. The first shell has the vector $(0,1,2,-2,-1)$, with
\[
s_2=10,\quad s_3=0,\quad s_4=34,\quad s_5=0.
\]
Both polynomials vanish there. Their theta weights are $20$ and $22$. The cusp spaces are respectively spanned by $\Delta E_4^2$ and $\Delta E_4E_6$, each with nonzero first coefficient. The zero first coefficient therefore kills the whole theta form. This is also the simplest valence argument for these two degrees.

The first retained family has fourteen corrections, with counts
\[
8:1,\quad10:1,\quad12:1,\quad14:2,\quad16:2,\quad18:3,\quad20:4.
\]
The transport continuation adds $22$ more: four at degree $22$, five at $24$, six at $26$ and seven at $28$. The total is $36$. These are explicit permitted corrections, not a proof that every universally null polynomial lies in their span.

\section{Moment identities that introduce other forms}
The harmonic polynomial $D(130s_4-47s_2^2)$ has weight-$16$ theta form. Its first coefficient is $-280$. Write $E_4\Delta=\sum a_{16}(n)q^n$. Then
\begin{equation}\label{eq:moment16}
13\sum_{v\in\Vn}W(v)s_4(v)
=470n^2\tau(n)-28a_{16}(n).
\end{equation}
Likewise $D(90s_2s_4-27s_2^3-320s_3^2)$ gives
\begin{align}\label{eq:moment18}
900n\sum W(v)s_4(v)-27000n^3\tau(n)
-320\sum W(v)s_3(v)^2
=3600a_{18}(n),
\end{align}
where $E_6\Delta=\sum a_{18}(n)q^n$ and the packet sums are on $\Vn$. These are exact identities. Their additional cusp coefficients must be controlled if the identities are used to infer nonvanishing. A new modular moment is not automatically a sign or residue certificate.

\chapter{What an exact sign certificate proves}\label{ch:certificates}
\section{The sufficient sign condition}
Take a finite rational combination of universally null corrections and set $F=1+\sum a_jQ_j$. If all real numbers $W(v)F(v)$ on a shell have one weak sign and at least one is nonzero, their sum is nonzero. By \eqref{eq:correctedtotal}, that sum is $\tau(n)$.

This is a valid certificate. Its verification requires the universal null identities, the exact rational coefficients, every packet in the shell and every claimed sign. A numerical optimizer may propose coefficients, but the certificate is established by exact rational checks. A failed optimizer does not establish infeasibility without a separate exact obstruction.

\section{Six packets reduced to two at \texorpdfstring{$n=7$}{n=7}}
Choose
\[
a=\frac{383}{1411344000000},\qquad
b=\frac{139}{61158240000000},\qquad
F=1+aQ_8+bQ_{10}.
\]
There are three $J$-orbits, each of size two. Their representatives and weights are
\begin{center}
\begin{tabular}{lrr}
\toprule
Representative $v$&$W(v)$&$F(v)$\\
\midrule
$(-5,-4,2,3,4)$&1176&0\\
$(-5,6,2,-2,-1)$&$-8624$&0\\
$(0,-4,7,-2,-1)$&$-924$&$299/33$\\
\bottomrule
\end{tabular}
\end{center}
For the three rows, the $(Q_8,Q_{10})$ values are
\begin{align*}
&(712800000,-525096000000),\\
&(-3790800000,12636000000),\\
&(36288000000,-786240000000).
\end{align*}
Substitution verifies the displayed $F$ values exactly. Four packets disappear, and the remaining two are negative. Therefore
\begin{equation}\label{eq:tau7cert}
\tau(7)=2(-924)\frac{299}{33}=-16744.
\end{equation}
The solve can be performed by requiring only the first two rows to vanish. It need not use a preassigned value of $\tau(7)$. The surviving value then supplies the coefficient, and the independent product engine checks it.

An integral form of the same certificate multiplies $F$ by $183474720000000$:
\[
183474720000000+49790Q_8+417Q_{10}.
\]
Its surviving value is $1662392160000000$. Clearing denominators is useful for a short exact checker, but does not change the sign argument.

\figplaceholder{orbit7}{The three original $J$-orbit masses at $n=7$, each containing two packets. The exact certificate in the text removes the positive orbit and one negative orbit, then rescales the remaining negative one. The unchanged total is certified nonzero by the displayed rational coefficients.}

\section{An exact obstruction in a specified smaller space}
At $n=31$, the seven retained corrections through degree $16$ do not admit a nontrivial common-sign certificate. The shell has sixteen $J$-orbits. Let $\epsilon_i=\sgn W(v_i)$ for their representatives. The retained exact dual has strictly positive rational entries $y_i$ satisfying
\[
\sum_i y_i=1,\qquad\sum_i y_i\epsilon_i=0,\qquad
\sum_i y_i\epsilon_iQ_j(v_i)=0
\]
for all seven correction columns. Consequently, for every coefficient vector,
\[
\sum_i y_i\epsilon_i\left(1+\sum_j a_jQ_j(v_i)\right)=0.
\]
If every corrected packet had one weak sign and one were nonzero, every $\epsilon_iF(v_i)$ would have that sign and one would be strict. A strictly positive weighted average could not be zero. This contradiction proves infeasibility in the specified seven-column span, for either sign.

The scope matters. The dual is not a theorem about every polynomial of degree at most $16$, unless completeness of the retained list in that larger class is separately proved. It is not a theorem about all universally null polynomials. It does not show that $\tau(31)$ vanishes. It describes one exact finite correction space.

\section{The obstruction can disappear when the space changes}
Adding the retained corrections at degrees $18$ and $20$ gives an exact strictly negative certificate at $31$. The previous dual does not annihilate the larger family. This is an instructive example of a removable obstruction. It cannot support the assertion that every obstruction at every prime will be removable.

The retained strict rational certificates cover the following primes:
\begin{center}
\begin{tabular}{rrrrrrr}
\toprule
$p$&2&3&5&7&11&13\\
Degree&8&8&8&10&8&12\\
Sign&$-$&$+$&$+$&$-$&$+$&$-$\\
\midrule
$p$&17&19&23&29&31&37\\
Degree&8&14&14&16&20&20\\
Sign&$-$&$+$&$+$&$+$&$-$&$-$\\
\midrule
$p$&47&59&&&&\\
Degree&14&20&&&&\\
Sign&$+$&$-$&&&&\\
\bottomrule
\end{tabular}
\end{center}
The degrees used are not claimed optimal. In particular, no exact degree-$18$ impossibility at $31$ is asserted. The fresh release replay verifies all fourteen stored rational certificates and all $300$ packet margins in their combined finite scope.

\figplaceholder{orbit31}{The sixteen original $J$-orbit masses at $31$. Eight are positive and eight negative. The coefficients sum to the nonzero total $-52843168$, but individual nonzero masses and a balance of sign counts do not force that total.}

\section{Completeness retains the total}
Suppose a shell has nonzero orbit masses $w_1,\ldots,w_N$ and total $T$. Even allow every rational correction vector $q$ with $\sum_iw_iq_i=0$. Then
\begin{equation}\label{eq:fullsignequiv}
\exists q:\ w_i(1+q_i)\text{ have one weak sign and are not all zero}
\quad\Longleftrightarrow\quad T\ne0.
\end{equation}
Necessity follows from conservation. For sufficiency, if $T\ne0$, take
\[
q_i=\frac{T}{Nw_i}-1.
\]
Every corrected mass is $T/N$ and the correction has zero weighted total. This construction proves an equivalence. If used as a proof of $T\ne0$ without independent information about $T$, it would be circular.

Thus a theorem that all zero-total correction vectors are available would explain the existence of a sign certificate once the total is known to survive. It would not independently establish that survival. A concrete sign certificate can still be decisive, as at $7$. A general completeness theorem must be combined with an independently established nonzero sign or arithmetic condition to settle every prime.
